\section*{Bab \ref{ch:g11:drift-to-spreading} --- Dari Pergeseran Benua sampai Pemekaran Dasar Samudra}

\begin{solution}{exo:g11:drift-to-spreading:1}
\omterm{def:g11:drift-to-spreading:spreading}{Pemekaran dasar samudra} adalah pembentukan dasar samudra baru di puncak
punggung tengah samudra, tempat batuan cair naik dan membeku, yang diikuti
oleh geraknya menjauhi punggung itu di kedua sisi. Laju penuh adalah laju dua
lempeng saling menjauh; laju setengah adalah laju dasar di satu sisi menjauhi
sumbu. Untuk pemekaran yang simetris, laju penuh $= 2
\times$ laju setengah.
\end{solution}

\begin{solution}{exo:g11:drift-to-spreading:2}
Di atas \omterm{def:g11:drift-to-spreading:remanence}{suhu Curie-nya}, sekitar \qty{580}{\celsius}, magnetit tidak dapat
mempertahankan magnetisasinya sendiri; hanya ketika mendingin di bawah suhu itu
butir-butirnya termagnetisasi searah medan dan mempertahankan arah itu. Bila
dipanaskan lagi sampai \qty{700}{\celsius}, di atas \omterm{def:g11:drift-to-spreading:remanence}{suhu Curie}, batuan itu
kehilangan rekaman lamanya; ketika mendingin lagi, batuan itu merekam medan
pada masa yang baru itu.
\end{solution}

\begin{solution}{exo:g11:drift-to-spreading:3}
Bunyi itu turun dan kembali: $d = 1500 \times 5.2/2 = \qty{3900}{m}$.
\end{solution}

\begin{solution}{exo:g11:drift-to-spreading:4}
Selang normal masa kini bermula \num{0.78} juta tahun yang lalu: selang itu
telah berlangsung \num{780000} tahun. Peristiwa Jaramillo berlangsung
$1.07 - 0.99 = \num{0.08}$ juta tahun, yaitu \num{80000} tahun. Dengan
menghitung batas-batas yang lebih muda daripada lima juta tahun (0.78, 0.99,
1.07, 1.77, 1.95, 2.14, 2.15, 2.581, 3.04, 3.11, 3.22, 3.33, 3.58, 4.18, 4.29,
4.48, 4.62, 4.80, 4.89, 4.98): dua puluh pembalikan.
\end{solution}

\begin{solution}{exo:g11:drift-to-spreading:5}
$65\,000/40\,000 \approx 1.6$: punggung-punggung itu akan mengelilingi Bumi
sekitar satu setengah kali.
\end{solution}

\begin{solution}{exo:g11:drift-to-spreading:6}
$20/0.78 \approx \qty{25.6}{km}$ per sejuta tahun, yaitu laju setengah sekitar
\qty{2.6}{cm} setahun; laju penuhnya sekitar \qty{5.1}{cm} setahun.
\end{solution}

\begin{solution}{exo:g11:drift-to-spreading:7}
Barat: $45/3.58 \approx 12.6$~km per sejuta tahun, \qty{1.26}{cm} setahun;
timur: $44/3.58 \approx 12.3$~km per sejuta tahun, \qty{1.23}{cm} setahun.
Laju penuh: $1.26 + 1.23 \approx \qty{2.5}{cm}$ setahun. Kedua laju setengah
itu berbeda sekitar \qty{2}{\percent}: pemekarannya hampir simetris.
\end{solution}

\begin{solution}{exo:g11:drift-to-spreading:8}
\qty{1.9}{cm} setahun sama dengan \qty{19}{km} per sejuta tahun:
$990/19 \approx 52$ juta tahun, sedikit lebih tua daripada 49 juta tahun umur
sedimen tertuanya (sekitar \qty{6}{\percent}).
\end{solution}

\begin{solution}{exo:g11:drift-to-spreading:9}
Dengan \qty{20}{km} per sejuta tahun: \qty{15.6}{km}, \qty{19.8}{km},
\qty{21.4}{km}, \qty{35.4}{km}, dan \qty{39.0}{km}. Dari sumbu: normal dari 0
sampai \qty{15.6}{km}, terbalik sampai \qty{19.8}{km}, normal sampai
\qty{21.4}{km} (Jaramillo), terbalik sampai \qty{35.4}{km}, normal sampai
\qty{39.0}{km} (Olduvai), lalu terbalik.
\end{solution}

\begin{solution}{exo:g11:drift-to-spreading:10}
$64 \times 10^6/4.6 \times 10^9 \approx \num{0.014}$, sekitar
\qty{1.4}{\percent} umur Bumi. Hal ini mendukung argumen Hess bahwa dasar
samudra itu muda karena selalu dibuat dan dihancurkan: tidak ada dasar samudra
yang bertahan dari sebagian besar sejarah Bumi.
\end{solution}

\begin{solution}{exo:g11:drift-to-spreading:11}
Pada setiap saat, dasar terbentuk di sumbu dengan polaritas saat itu, lalu
terbelah: separuhnya bergerak menjauh di setiap sisi dengan laju yang sama.
Karena itu dua jalur yang terbentuk pada saat yang sama terletak pada jarak
yang sama dari sumbu, satu di setiap sisi, dengan polaritas yang sama: setiap
sisi adalah salinan sisi lainnya, yang dicerminkan pada sumbu. Simetri itu
rusak jika dasar bergerak menjauh lebih cepat di satu sisi daripada di sisi
lain (pemekaran asimetris), atau jika sumbunya melompat ke kedudukan baru.
\end{solution}

\begin{solution}{exo:g11:drift-to-spreading:12}
(1) \qty{5}{cm} sama dengan $5 \times 10^{-5}$~\unit{km}: setiap tahun
ditambahkan jalur seluas $65\,000 \times 5 \times 10^{-5} = \qty{3.25}{km^2}$.
(2) $2.9 \times 10^8/3.25 \approx 9 \times 10^7$ tahun, sekitar 90 juta
tahun. (3) Ini sama ordenya dengan umur dasar tertua (190 juta tahun) dan
lebih singkat daripadanya: agar dasar samudra tetap begitu muda, dasar yang
dihancurkan, di suatu tempat, harus sebanyak yang dibuat di punggung-punggung.
\end{solution}

\begin{solution}{exo:g11:drift-to-spreading:13}
Peristiwa itu berlangsung \num{0.08} juta tahun. Dengan \qty{12.5}{km} per
sejuta tahun garisnya selebar $0.08 \times 12.5 = \qty{1}{km}$; dengan
\qty{75}{km} per sejuta tahun, $0.08 \times 75 = \qty{6}{km}$. Dilihat dari
permukaan laut, \qty{3}{km} di atasnya, garis yang jauh lebih sempit daripada
kedalamannya hanya menghasilkan medan yang lemah, yang mengabur bersama medan
garis-garis tetangganya yang lebar dan berpolaritas berlawanan; garis yang
dua kali lebih lebar daripada kedalamannya tampak menonjol.
\end{solution}

\begin{solution}{exo:g11:drift-to-spreading:14}
Lokasi 16: $176/11 \approx \qty{16}{m}$ per sejuta tahun; lokasi 20: $72/67
\approx \qty{1.1}{m}$ per sejuta tahun. Sedimen menumpuk dengan laju yang
sangat berbeda dari tempat ke tempat dan dari waktu ke waktu, sehingga tebalnya
tidak memberikan umur dasar samudra. Umur itu diberikan oleh fosil-fosil
sedimen tertua, yang terletak langsung di atas basalt.
\end{solution}

\begin{solution}{exo:g11:drift-to-spreading:15}
(1) Barat: $58/2.581 \approx \qty{22.5}{km}$ per sejuta tahun,
\qty{2.25}{cm} setahun; timur: $70/2.581 \approx \qty{27.1}{km}$ per sejuta
tahun, \qty{2.71}{cm} setahun; laju penuh $128/2.581 \approx \qty{49.6}{km}$
per sejuta tahun, sekitar \qty{5.0}{cm} setahun. (2) Sekitar \qty{225}{km} di
barat dan \qty{271}{km} di timur. (3) Kedua batas itu berimpit di sumbu ketika
pembalikan terjadi; sejak itu keduanya saling menjauh dengan laju penuh, bagaimanapun
gerak itu terbagi di antara kedua sisi. Jarak pisah keduanya dibagi umurnya
adalah laju penuh.
\end{solution}

\begin{solution}{pb:g11:drift-to-spreading:1}
\textbf{1.} Lokasi 16 yang paling dekat (\qty{191}{km}, 11 juta tahun), lokasi
20 yang paling jauh (\qty{1270}{km}, 67 juta tahun): makin jauh dari punggung,
makin tua dasarnya.

\textbf{2.} Sedimen mulai mengendap di atas basalt segera setelah basalt itu
terbentuk, sehingga lapisan tertua kira-kira setua basalt itu. Lapisan itu
hanya dapat sedikit lebih muda, tidak pernah lebih tua.

\textbf{3.} Dalam kilometer per sejuta tahun: lokasi 16: $191/11 \approx 17.4$;
lokasi 15: $380/24 \approx 15.8$; lokasi 18: $506/26 \approx 19.5$; lokasi 17:
$643/33 \approx 19.5$; lokasi 14: $727/40 \approx 18.2$; lokasi 19: $990/49
\approx 20.2$; lokasi 20: $1270/67 \approx 19.0$.

\textbf{4.} $129.5/7 \approx \num{18.5}$~km per sejuta tahun, sekitar
\qty{1.85}{cm} setahun.

\textbf{5.} Ya: perbandingannya (19.5) terletak di antara perbandingan
lokasi-lokasi barat. Dasar samudra telah memekar dengan laju yang kira-kira
sama di kedua sisi: pemekarannya simetris.

\textbf{6.} $t = d/u$: umurnya sebanding dengan jaraknya.

\textbf{7.} $u = 208\,276/10\,952 \approx \num{19.0}$~km per sejuta tahun,
yaitu \qty{1.90}{cm} setahun.

\textbf{8.} Sekitar $2 \times 1.9 = \qty{3.8}{cm}$ setahun.

\textbf{9.} Ya: \qty{1.9}{cm} setahun sedikit kurang dari \qty{2}{cm}
setahun.

\textbf{10.} $1270/19.0 \approx 67$ juta tahun, sama dengan umur sedimen
tertuanya.

\textbf{11.} $191/9 \approx \num{21.2}$~km per sejuta tahun,
\qty{2.1}{cm} setahun.

\textbf{12.} $1270/71 \approx \num{17.9}$~km per sejuta tahun,
\qty{1.8}{cm} setahun.

\textbf{13.} $180\,329/9854 \approx \num{18.3}$~km per sejuta tahun,
\qty{1.83}{cm} setahun.

\textbf{14.} $(1.90 - 1.83)/1.83 \approx \qty{4}{\percent}$: kedua laju itu
sesuai dalam batas ketidakpastian umur dan jarak.

\textbf{15.} Satu umur berasal dari fosil-fosil dalam sedimen, yang lain dari
kemagnetan basalt dan skala waktu pembalikan. Kesalahan dalam satu metode
tidak akan muncul dalam metode lainnya: kesesuaian keduanya adalah uji yang
kuat bagi \omterm{def:g11:drift-to-spreading:spreading}{pemekaran dasar samudra}.

\textbf{16.} $13.8 + 13.8 = \ang{27.6}$.

\textbf{17.} $111.2 \times \cos 30^\circ \approx \qty{96.3}{km}$ per derajat;
$27.6 \times 96.3 \approx \qty{2660}{km}$.

\textbf{18.} $2660/19.0 \approx 140$ juta tahun.

\textbf{19.} $2660/130 \approx \num{20.5}$~km per sejuta tahun, sekitar
\qty{2.0}{cm} setahun, sedikit lebih besar daripada \qty{1.9}{cm} setahun.
Karena 67 juta tahun terakhir berlangsung dengan laju sekitar \qty{1.9}{cm}
setahun, pembukaan awal itu sedikit lebih cepat: dasar di antara lokasi 20 dan
tepi samudra, $2660 - 1270 =
\qty{1390}{km}$, terbentuk dalam $130 - 67 = 63$ juta tahun, sekitar
\qty{2.2}{cm} setahun.

\textbf{20.} \textbf{Dasar samudra tertua di dekat Afrika pada \ang{30} LS
berumur sekitar 130 juta tahun} (140 juta tahun bila laju masa kini
diekstrapolasi): Atlantik Selatan mulai membuka di sana pada periode Kapur.
\end{solution}
